\section{Introduction}
\label{sec:intro}

% ~0.5 page. Notes from OUTLINE.md; turn them into prose.
\todo{Paragraph 1, why travel time on foot matters: search and rescue,
wildfire escape routes, trail apps, accessibility maps. GIS still uses
Tobler's 1993 hiking function~\cite{tobler1993},
$v = 6\exp(-3.5\,|s + 0.05|)$ km/h, and fixed multipliers for other
modes.}

\todo{Paragraph 2, what is known: crowdsourced GPS shows Tobler is too fast
and too peaked~\cite{campbell2019crowdsourced,campbell2022variability};
Wood et al.~\cite{wood2023hiking} add OSM terrain classes and lidar
obstruction (England and Wales); Huang et al.~\cite{huang2026modeling} fit
boosted trees with 16 factors, weather included (Hong Kong); \v{S}eri\'c et
al.~\cite{seric2022modeling} predict a latent walking speed from Sentinel-2
(one Croatian range); HikingTTE~\cite{asako2024hikingtte} learns whole-hike
time from the hiker's first part. STRIDE~\cite{campbell2024stride} models
walking speed from lidar slope, understory
density~\cite{campbell2018understory} and roughness, calibrated on field
walks; RasterNet~\cite{hadzic2020rasternet} learns road speed from lidar
and imagery.}

\todo{Paragraph 3, the gap: each of these studies fits and tests in one
country, city or range with random or within-area splits (Wood: 10-fold
cross-validation; Huang: random split); none tests whether terrain
covariates transfer to regions a model has never seen, against the map tags
used today, with calibrated uncertainty. Wood et al. also found no gain over
existing functions for whole-route times (13.5--15.5\% error for every
model).}

\todo{Paragraph 4, contributions (three bullets): (1) a 17-region dataset of
walking traversals of 50~m trail segments from public OSM traces; (2) an
evaluation protocol (leave one region out, CRPS, coverage) with published
baselines; (3) the finding that elevation-derived terrain transfers and map
tags do not, plus a one-input terrain-adjusted hiking function.}
